\songtitle{Los rebotes}

\tag*{2026}\tag{ricochets}
\begin{notes}
Trova antioqueña summary by Claude with instructions, refinmentes, and some manual editions by Carlos Thompson %
summarizing the first 12 chapters of \emph{Ricochets}.
\end{notes}
\begin{style}
Trova antioqueña, Acoustic guitar plays rhythmic strumming patterns with frequent percussive slaps on the strings, Two male vocalist sing in a clear, narrative style with moderate vibrato, alternating verses, The tempo is 125 BPM in 4/4 time, The arrangement features a call-and-response dynamic between the lead vocals
\end{style}
\begin{lyrics}
\annotation{I}\\
John despierta entre el guayabo,\\
tirado sobre el piso frío;\\
Camilo lo alza del suelo,\\
sin juzgar su desvarío.

\annotation{II}\\
Sin juzgar su desvarío\\
guarda el gringo su tristeza;\\
la botella quedó vacía,\\
igual quedó su cabeza.

\annotation{III}\\
Igual quedó su cabeza\\
recordar a Lilith no alcanza;\\
en la pieza de Candelaria\\
el pasado no lo deja en calma.

\annotation{IV}\\
El pasado no lo deja en calma,\\
llega Lida con su psicología;\\
le pide consejo el soldado,\\
y ella acepta con alegría.

\annotation{V}\\
Y ella acepta con alegría\\
oír la historia de su condena:\\
la esposa que murió en Miami,\\
la niña que cayó por su pena.

\annotation{VI}\\
La niña que cayó por su pena\\
tiene nombre: Sarah, la abogada;\\
Harris guarda el mismo dolor,\\
frente a la carpeta encontrada.

\annotation{VII}\\
Frente a la carpeta encontrada\\
de MCT y su lavado,\\
Harris recuerda los tiempos\\
en que a Sarah había amado.

\annotation{VIII}\\
En que a Sarah había amado,\\
Müller en su hotel de Nueva York\\
con Sophie Durand conversa,\\
mientras negocia su valor.

\annotation{IX}\\
Mientras negocia su valor\\
Little pregunta por la noche\\
en que el tiroteo de Miami\\
dejó tres muertos en su broche.

\annotation{X}\\
Dejó tres muertos en su broche\\
Gunther, la niña y su pistola;\\
Keaton y Gonzalez indagan\\
quién fue el que armó esa bola.

\annotation{XI}\\
Quién fue el que armó esa bola\\
pregunta Perez al ya preso;\\
Müller responde con cautela,\\
sin traicionar ningún proceso.

\annotation{XII}\\
Sin traicionar ningún proceso\\
John en Bogotá lo sigue,\\
Jairo le cuenta los planes\\
de a quién Müller recibe.

\annotation{XIII}\\
De a quién Müller recibe\\
entra a la casa el vigilante,\\
copia el correo con Perez\\
y sale sin dejar instante.

\annotation{XIV}\\
Y sale sin dejar instante\\
Gebhart le habla con cautela,\\
de MCT algo revela\\
y entre la gente se vuela.

\annotation{XV}\\
Y entre la gente se vuela,\\
llega toda la familia entera;\\
don Jorge y doña Antonina\\
llegan con fe que persevera.

\annotation{XVI}\\
Llegan con fe que persevera,\\
Carolina llega de Alemania,\\
cruza el aire hasta Bogotá\\
con su acento y su nostalgia.

\annotation{XVII}\\
Con su acento y su nostalgia,\\
John y Lida van al cine;\\
juntos rompen la rutina\\
mientras el amor entre ellos define.

\annotation{XVIII}\\
Mientras el amor entre ellos define\\
Müller y Little hablan del negocio,\\
de Inter-Flota y sus caminos\\
entre el miedo y el reposo.

\annotation{XIX}\\
Entre el miedo y el reposo,\\
Carolina pregunta y espera;\\
sobre el papel él le confiesa\\
que hay micrófonos que acechan.

\annotation{XX}\\
Que hay micrófonos que acechan,\\
sabe Lida al partir un día;\\
se despide de John en la puerta\\
y el amor queda en la memoria fría.
\end{lyrics}

